% =====================================================================
%  ANNEXES — A : formats complets · B : glossaire
% =====================================================================
\appendix
\section{Les formats, octet par octet}\label{sec:formats}

Ces listings sont la reprise \emph{exacte} des formats normatifs des specs scellées
(\code{v06/I/CF2/README.md} § 3--4, \code{v06/I/CF3/README.md} § 3--4) — reproduits
ici parce que le maître les recopie, les compare et les fait lire à un tiers. En cas
de doute : l'arbre signé fait foi.

\subsection{Conteneur TRIADE-CF1}

\begin{lstlisting}
magic "TRIADE-CF1" . version u16 . champs longueur-prefixee u16 :
    hachage (sha256), kdf (argon2id + hkdf-sha256),
    chiffrement (xchacha20poly1305), date ISO-8601 UTC .
sel argon 16/32 o (pose a l'init) . params (t, m, p) .
sel HKDF 32 o . nonce 24 o (frais a CHAQUE ecriture) .
empreinte 32 o = sha256(chiffretexte) .
chiffretexte = AEAD(clair, K, AAD = tout ce qui precede)
    K = HKDF-SHA256(Argon2id(phrase, sel, 65536/3/1),
                    sel_hkdf, info "coffre=v1")
\end{lstlisting}

\subsection{Conteneur TRIADE-CF2}

\begin{lstlisting}
magic "TRIADE-CF2" (10 o) . version u16 (1) .
champs longueur-prefixee u16 : hachage (sha256), kdf (hkdf-sha256),
                               chiffrement (xchacha20poly1305),
                               date ISO-8601 UTC, shamir ("t=3/n=5") .
sel HKDF 32 o (aleatoire a l'init) . nonce 24 o (frais a CHAQUE
ecriture) . empreinte sha256 du chiffretexte 32 o .
verif 16 o = sha256(HKDF(M, sel, info "cf2=verif")) tronque .
chiffretexte : XChaCha20-Poly1305(JSON base, K, AAD = tout ce qui precede)
   avec K = HKDF-SHA256(M, sel du conteneur, info "coffre=v2", L=32)
\end{lstlisting}

\subsection{Fragment TRIADE-CF2-FRAG (ASCII, transportable sur papier)}

\begin{lstlisting}
TRIADE-CF2-FRAG v1
coffre: <sel HKDF du conteneur, 64 hex - l'identite STABLE du coffre>
x: <indice, decimal>
seuil: t=3/n=5
fragment: <64 hex - 32 o>
tag: <16 hex - sha256("TRIADE-CF2-frag" || x || fragment) tronque>
\end{lstlisting}

\subsection{Conteneur TRIADE-CF3 (petit-boutiste)}

\begin{lstlisting}
"TRIADE-CF3" 10 o | version u16 (=1) | epoque u64 | H_projet 32 o |
sel_argon 16 o | params (t,m,p : u32x3) | sel_hkdf 32 o | S_serrure 32 o |
nonce 24 o |                       <- fin de l'AAD (168 o)
empreinte 32 o = sha256(chiffretexte) | verif 16 o |
chiffretexte = AEAD(clair, aad, nonce, KEK)

KEK   = HKDF-SHA256(ikm = K_user || S_serrure || S_piege ;
                    salt = sel_hkdf ;
                    info = "triade-CF3/v1|" || epoque(le64) || H_projet)
verif = sha256(HKDF(ikm = K_user || V_piege ;
                    info = "triade-CF3/verif/v1"))[:16]
\end{lstlisting}

\subsection{Piège TRIADE-CF3-PIEGE (ASCII, un fichier par époque)}

\begin{lstlisting}
TRIADE-CF3-PIEGE v1
epoque: <u64 decimal>
projet: <H_projet hex (64)>
piege: <S_piege hex (64)>
verif: <V_piege hex (32)>
tag: <hex (16)>
tag = sha256("TRIADE-CF3-piege" || epoque(le64) || H_projet ||
             S_piege || V_piege)[:8]
\end{lstlisting}

\subsection{Manifeste du sceau de l'Intention (canonique)}

\begin{lstlisting}
TRIADE-SCEAU-I v1
arbre: <nom>
artefact: <artefact>
balise: <balise>
fichiers: <N>
sha256:<64 hex minuscules>  <chemin relatif POSIX, tri lexicographique>
...   (une ligne par fichier - contenu canonique LF)
\end{lstlisting}

L'empreinte d'arbre est le sha256 du manifeste \emph{entier} ; la signature Ed25519
(RFC 8032) couvre l'empreinte — 32 octets, imprimables.

\section{Glossaire du coffre}\label{sec:glossaire}

\begin{description}[style=nextline, leftmargin=1.2em]
\item[acte] Une opération qui déverrouille (ouvre, clôture, vérifie) — porte la
sentinelle instantanée.
\item[carte du sachet] Fiche papier du sachet où le maître consigne, \emph{à la main},
le checksum de 8 hex de chaque clôture — deuxième témoin d'époque, indécevable par
logiciel.
\item[clé I] La graine Ed25519 hors-ligne qui signe l'intention (la spécification) ;
le dépôt ne voit que sa clé publique.
\item[époque] Compteur u64 du coffre CF3, posé à 1, incrémenté à chaque clôture —
l'horloge de la métamorphose et le témoin anti-retour.
\item[fail-closed] Tout refus est nommé et rien ne s'écrit : l'échec est sûr, jamais
silencieux.
\item[fragment] Un morceau ASCII du secret maître CF2 (32 octets utiles + identité) ;
transportable sur papier.
\item[garde] Répertoire distinct portant UN fragment CF2 (mode 0700) ; en production,
ACL et supports séparés.
\item[journal local] \code{.keys/journal-cf3.txt} — premier témoin d'époque, contrôle
automatique, effaçable par un attaquant admin.
\item[maison] Le clone git (avec \code{v06/} entier) ou le kit Drive — l'endroit d'où
les cérémonies se lancent.
\item[média du piège] Le support séparé (clef USB du sachet) qui porte S\_piège et
V\_piège — jamais écrit sur la machine hors l'instant de l'acte.
\item[piège] Le fichier \code{piege-CF3-epoque-N.txt} du média — sans lui,
l'acte est impossible.
\item[PV] Procès-verbal de destruction : date, déclencheur, artefacts détruits,
raison nommée — deux destructions même seconde, deux fichiers (D123).
\item[relève] Ce qui ressuscite le coffre si la machine meurt : le sachet (pièges,
copies hors-ligne) et le fichier-remise D81.
\item[sachet] La pochette physique individuelle : pièges d'époque, copies hors-ligne
des fragments, la carte.
\item[sceau de l'Intention] La signature Ed25519 des arbres de couche I ; dix arbres,
gardes de build et de patch.
\item[sentinelle instantanée] Les cinq vérifications ordonnées qui courent à chaque
acte CF3 — par invocation, jamais en démon.
\item[seuil] Le T de Shamir : le nombre de fragments qui reconstituent M (et le nombre
de copies hors-ligne qui arment la relève).
\item[staging] Le conteneur en attente (\code{coffre\_cf3.bin.staging}) pendant la
re-lecture à froid ; remplacé atomiquement ou détruit.
\item[témoins d'époque] Les trois miroirs anti-retour : journal local, carte papier,
journal Git.
\item[trio D122] Une GUI, une boîte noire, un bouton « Envoyer le rapport » — pour
chaque module à tester.
\item[vérités vivantes] Les modules CLI que les GUI relaient sans réimplémenter
(\code{coffre\_cf2.py}, \code{coffre\_cf3.py}, \code{sceau\_I.py}).
\end{description}

\section{Les questions du terrain — et leurs réponses}\label{sec:faq}

Posées par le terrain réel entre le 30 et le 31 octobre 2026, répondues alors au fil
de l'eau — écrites ici une fois pour toutes. Les mêmes questions reviendront ; les
mêmes réponses tiennent.

\begin{description}[style=nextline, leftmargin=1.2em]
\item[« La cérémonie m'a produit un \code{fragments-cf2.txt} et le système attend
3 \code{fragment-*.cf2} — où sont les fragments ? »] Les CINQ originaux vivent sous
\code{.keys\textbackslash gardes-cf2\textbackslash g1...g5\textbackslash
fragment-N.cf2f} — un fichier par garde, dossier caché. La feuille est l'export
\emph{papier} des trois fragments g3--g5 ; les copies USB sont les fichiers
\code{.cf2f} eux-mêmes. Trois objets, trois rôles (§ \ref{sec:cf2}).
\item[« Sur la page de choix de scénario, il y a deux boutons d'attachement qui
font doublon. »] Vu et corrigé : depuis D131, un bouton PAR page — l'accueil ne le
porte que dans le panneau de scénario, la page de cérémonie que dans son bandeau.
Le geste, lui, n'a pas changé.
\item[« L'interface de migration n'est pas claire ; les boutons de recherche ne sont
pas alignés aux champs. »] Défaut mécanique du v1.0 (champs pleine largeur, boutons
rejetés en colonne à droite) — corrigé par la grille D132 : chaque bouton sur la
ligne de son champ, confirmation de phrase étiquetée, actes pleine largeur. La
fenêtre s'annonce v1.1.
\item[« \code{lancer\_migration\_cf3.bat} n'est pas dans git. »] Il y est — à la
\emph{racine} du dépôt, à côté de \code{v06\textbackslash}, pas DANS
\code{v06\textbackslash} : convention D118. \code{git ls-files} le liste ;
\code{git log} le suit depuis D127.
\item[« La mise à jour est un \code{.py} que le \code{.bat} n'appelle pas. »] Le
lanceur appelle \code{v06\textbackslash migration\_cf3\_gui.py} \emph{en premier}
(ligne 12 du \code{.bat}) — et le correctif vivait précisément dans ce fichier.
Vérifiable en trois endroits à l'exécution : titre, bandeau, première ligne de la
boîte noire.
\item[« Il demande un \code{coffre\_cf3.bin} que je n'ai pas. »] Ce champ est une
\emph{sortie} : « le coffre qui naît », dialogue d'enregistrement, pré-rempli. Ce
qui manquait vraiment dans ce rapport-là était \code{coffre\_cf3.py} — le module :
le GUI tournait hors de sa maison (§ \ref{sec:maison}).
\item[« Modules REQUIS absents : \code{coffre\_cf3.py}, \code{coffre\_cf2.py}. »]
Le GUI est un relais : sans ses vérités vivantes à côté de lui, chaque acte est un
refus nommé. Remède : lancer depuis la racine du clone (ou du kit), où
\code{v06\textbackslash} est entier.
\item[« Je dois pousser quelque chose sur git maintenant ? »] Non — par
construction : tout ce que la migration écrit vit dans \code{.keys\textbackslash}
(gitignoré) ou sur le média ; la boîte noire est un \code{*.log} (ignoré).
\code{git status} rend « nothing to commit ». Les seuls push de votre chaîne sont
les remises du sceau.
\item[« C'est quoi, la carte du sachet ? »] Une fiche papier, dans le sachet, où
vous consignez À LA MAIN le checksum de 8 hex de chaque clôture — le témoin
qu'aucun logiciel ne peut truquer (§ \ref{sec:gards}).
\item[« On n'avait pas décidé de supprimer la sentinelle en migrant en CF3 ? Si on
impose CF3 dès la création, on n'a plus besoin de sentinelle, non ? »] Le
\emph{démon} est mort en CF3 ; l'exigence D120 vit \emph{par invocation} : la
sentinelle instantanée est la serrure — impossible à oublier de lancer. Et oui :
à l'époque de cette question, un coffre neuf naissait CF3 par \code{init} direct ;
depuis D137 (v0.6.1), le niveau de naissance — CF1, CF2 ou CF3 — est votre
arbitrage à la création (spécification § 12.5).
\item[« J'ai fermé sans exporter le rapport, mais ça a marché. »] Normal : l'export
n'est pas un prérequis de la cérémonie — la vérité vit dans le conteneur, les
gardes et le journal local. Le rapport est la \emph{mémoire} du geste : le bouton
existe pour que le prochain imprévu ait son dossier complet.
\item[« Un premier \code{migrer} a échoué : phrase de contrôle divergente. »] Le
fail-closed qui travaille : la remise D81 n'était pas encore remplacée. Rien n'a été
écrit ; posez la phrase courante dans le fichier-remise, commit, push — puis
rejouez.
\end{description}
